% Môn: Toán
% Lớp: 11
% Chương: Chương I. Hàm số lượng giác và phương trình lượng giác
% Bài: Bài 3. Hàm số lượng giác
% Số câu: 17
% App đọc file này trực tiếp — sửa rồi Commit trên GitHub, bấm Đồng bộ GitHub .tex

% Bài 3. Hàm số lượng giác

% ===== Câu 1 =====
% ID: T11C1B3-01-DS
% Mức: VD
\dangbt{Sử dụng công thức cộng và trừ của tang}
\begin{ex}
% ID: T11C1B3-01-DS
% Mức: VD
Cho biết $\sin \alpha=\dfrac{3}{5}$, $\dfrac{\pi}{2}\lt \alpha<\pi$. Khi đó
\choiceTF

\loigiai{
A. ch Đúng. Vì $\dfrac{\pi}{2}\lt \alpha \lt \pi$ nên $\cos \alpha \lt 0$.
B. ch Đúng. Ta có $\cos \alpha =-\sqrt{1-\sin^2\alpha}=-\dfrac{4}{5}$.
C. ch Sai. Ta có $\tan \alpha =\dfrac{\sin \alpha}{\cos \alpha}=-\dfrac{3}{4}$.
D. ch Sai. Ta có $\tan \left(\alpha+\dfrac{\pi}{3}\right)=\dfrac{\tan \alpha+\tan \dfrac{\pi}{3}}{1-\tan \alpha \tan \dfrac{\pi}{3}}=\dfrac{\tan \alpha+\sqrt{3}}{1-\sqrt{3} \tan \alpha}$.
.
}
\end{ex}

% ===== Câu 2 =====
% ID: T11C1B3-01-TN
% Mức: NB
\dangbt{Sử dụng công thức lượng giác trong tam giác}
\begin{ex}
% ID: T11C1B3-01-TN
% Mức: NB
Biết $A$, $B$, $C$ là các góc của tam giác $ABC$, khi đó.
\choice
{$\sin C=-\sin (A+B)$}
{$\cos C=\cos (A+B)$}
{$\tan C=\tan (A+B)$}
{\True $\cot C=-\cot (A+B)$}
\loigiai{
Vì $A$, $B$, $C$ là các góc của tam giác $ABC$ nên $A+B+C=180^o\Rightarrow C=180^o-(A+B)$.
Do đó $C$ và $(A+B)$ là 2 góc bù nhau.
$\Rightarrow \sin C=\sin (A+B)$; $\cos C=-\cos (A+B)$; $\tan C=-\tan (A+B)$; $\cot C=-\cot (A+B)$.
}
\end{ex}

% ===== Câu 3 =====
% ID: T11C1B3-02-DS
% Mức: VD
\dangbt{Sử dụng công thức cộng và trừ của tang}
\begin{ex}
% ID: T11C1B3-02-DS
% Mức: VD
Cho $\cot x=-\sqrt{3}$, $\dfrac{3 \pi}{2}\lt x<2 \pi$. Khi đó
\choiceTF

\loigiai{
A. ch Đúng. $\cot x=-\sqrt{3}$, $\dfrac{3\pi}{2}\lt x<2\pi$. Khi đó $\dfrac{1}{\sin^2x}=1+\cot^2x=1+(-\sqrt{3})^2=10\\Rightarrow \sin^2x=\dfrac{1}{10}\Rightarrow \sin x=\pm \dfrac{\sqrt{10}}{10}$.
.
B. ch Sai. $\cos x=\cot x\cdot \sin x=-\sqrt{3}\cdot \left(-\dfrac{\sqrt{10}}{10}\right)=\dfrac{\sqrt{30}}{10}$.
C. ch Đúng.
D. ch Đúng.
}
\end{ex}

% ===== Câu 4 =====
% ID: T11C1B3-02-TN
% Mức: VDC
\dangbt{Sử dụng công thức lượng giác trong tam giác}
\begin{ex}
% ID: T11C1B3-02-TN
% Mức: VDC
Biểu thức $\dfrac{3-4\cos 2\alpha +\cos4\alpha}{3+4\cos 2\alpha +\cos4\alpha}$ có kết quả rút gọn bằng
\choice
{$-\tan^4\alpha$}
{\True $\tan^4\alpha$}
{$-\cot^4\alpha$}
{$\cot^4\alpha$}
\loigiai{
Ta có 
 \begin{eqnarray*}
 && $\dfrac{3-4\cos 2\alpha +\cos4\alpha}{3+4\cos 2\alpha +\cos4\alpha}=\dfrac{3-4(1-2\sin^2\alpha)+2(1-2\sin^2\alpha)^2-1}{3+4(2\cos^2\alpha -1)+2(2\cos^2\alpha -1)^2-1}& =&\dfrac{8\sin^2a-8\sin^2\alpha +8\sin^4\alpha}{8\cos^2a-8\cos^2\alpha +8\cos^4\alpha}=\tan^4\alpha$. 
 \end{eqnarray*}
}
\end{ex}

% ===== Câu 5 =====
% ID: T11C1B3-03-DS
% Mức: VDC
\dangbt{Sử dụng công thức cộng và trừ của tang}
\begin{ex}
% ID: T11C1B3-03-DS
% Mức: VDC
Biết $\sin a=\dfrac{8}{17}$, $\tan b=\dfrac{5}{12}$ và $a$, $b$ là các góc nhọn. Khi đó
\choiceTF

\loigiai{
A. ch Đúng. Vì $a$, $b$ là các góc nhọn nên $\cos a>0$, $\cos b>0$.
.
B. ch Đúng. $\cos b=\sqrt{\dfrac{1}{1+\tan ^2 b}}=\dfrac{12}{13} \Rightarrow \sin b=\cos b \tan b=\dfrac{5}{13}$.
Khi đó $\sin (a-b)=\sin a \cos b-\cos a \sin b=\dfrac{8}{17} \cdot \dfrac{12}{13}-\dfrac{15}{17} \cdot \dfrac{5}{13}=\dfrac{21}{221}$.
C. ch Sai. $\cos (a+b)=\cos a \cos b-\sin a \sin b=\dfrac{15}{17} \cdot \dfrac{12}{13}-\dfrac{8}{17} \cdot \dfrac{5}{13}=\dfrac{140}{221}$.
D. ch Sai. $\tan (a+b)=\dfrac{\tan a+\tan b}{1-\tan a \tan b}=\dfrac{\dfrac{8}{15}+\dfrac{5}{12}}{1-\dfrac{8}{15} \cdot \dfrac{5}{12}}=\dfrac{171}{140}$.
}
\end{ex}

% ===== Câu 6 =====
% ID: T11C1B3-03-TN
% Mức: VDC
\dangbt{Sử dụng công thức lượng giác trong tam giác}
\begin{ex}
% ID: T11C1B3-03-TN
% Mức: VDC
Cho biểu thức $A=\sin^2(a+b)-\sin^2a-\sin^2b$. Hãy chọn kết quả đúng
\choice
{$A=2\cos a\cdot \sin b\cdot \sin (a+b)$}
{$A=2\sin a\cdot \cos b\cdot \cos (a+b)$}
{$A=2\cos a\cdot \cos b\cdot \cos (a+b)$}
{\True $A=2\sin a\cdot \sin b\cdot \cos (a+b)$}
\loigiai{
Ta có
 \begin{eqnarray*}
 A&=& $\sin^2(a+b)-\sin^2a-\sin^2b=\sin^2(a+b)-\dfrac{1-\cos 2a}{2}-\dfrac{1-\cos 2b}{2}&=&\sin^2(a+b)-1+\dfrac{1}{2}(\cos2\,\text{a}$ + $\cos$ 2b)=- $\cos^2(a+b)+\cos$ (a+b) $\cos$ (a-b)
&=& $\cos$ (a+b)$$[$\cos$(a-b)-$\cos$(a+b)$$]=2$\sin$a$\sin$b$\cos$ (a+b).
 \end{eqnarray*}
}
\end{ex}

% ===== Câu 7 =====
% ID: T11C1B3-04-DS
% Mức: VDC
\begin{ex}
% ID: T11C1B3-04-DS
% Mức: VDC
Biết $\cos 2 \alpha=\dfrac{5}{9}$, $0^\circ<\alpha<90^\circ$. Khi đó
\choiceTF
{\True $\sin \alpha=\dfrac{\sqrt{28}}{9}$}
{\True $\cos \alpha=\dfrac{\sqrt{53}}{9}$}
{$\tan \alpha =\dfrac{\sqrt{371}}{53}$}
{\True $\cot \alpha =\dfrac{\sqrt{371}}{14}$}
\loigiai{
\begin{itemchoice}
\item (Đúng) $\sin \alpha=\dfrac{\sqrt{28}}{9}$. (Vì): ch Đúng. Ta có $\cos 2 \alpha=\dfrac{5}{9}$, $0^\circ<\alpha<90^\circ$ nên
\item (Đúng) $\cos \alpha=\dfrac{\sqrt{53}}{9}$. (Vì): ch Đúng. $\cos ^2 \alpha=1-\sin ^2 \alpha=1-\left(\dfrac{\sqrt{28}}{9}\right)^2=\dfrac{53}{81} \Rightarrow \cos \alpha= \pm \dfrac{\sqrt{53}}{81}$
\item (Sai) $\tan \alpha =\dfrac{\sqrt{371}}{53}$. (Vì): ch Sai. $\tan \alpha=\dfrac{\sin \alpha}{\cos \alpha}=\dfrac{2 \sqrt{371}}{53}$
\item (Đúng) $\cot \alpha =\dfrac{\sqrt{371}}{14}$. (Vì): ch Đúng. $\cot \alpha=\dfrac{1}{\tan \alpha}=\dfrac{\sqrt{371}}{14}$
\end{itemchoice}
}
\end{ex}

% ===== Câu 8 =====
% ID: T11C1B3-04-TLN
% Mức: VDC
\begin{ex}
% ID: T11C1B3-04-TLN
% Mức: VDC
Trong Vật lí, phương trình tổng quát của một vật dao động điều hoà cho bởi công thức $x(t)=A \cos (\omega t+\varphi)$, trong đó $t$ là thời điểm (tính bằng giây), $x(t)$ là li độ của vật tại thời điểm $t$, $A$ là biên độ dao động $(A>0)$ và $\varphi \in[-\pi; \pi]$ là pha ban đầu của dao động. Xét hai dao động điều hoà có phương trình: $x_1(t)=3 \cdot \cos \left(\dfrac{\pi}{6} t+\dfrac{\pi}{6}\right)(cm)$, $x_2(t)=3 \cdot \cos \left(\dfrac{\pi}{6} t+\dfrac{\pi}{4}\right)(cm)$. Pha ban đầu của dao động tổng hợp $x(t)=x_1(t)+x_2(t)$ là (đơn vị độ)
\shortans{$37{,}5$}
\loigiai{
Ta có
 \begin{eqnarray*}
  x(t)&=&x_1(t)+x_2(t)=3 $\cos\left(\dfrac{\pi}{6}t+\dfrac{\pi}{6}\right)+3\cos\left(\dfrac{\pi}{6}t+\dfrac{\pi}{4}\right)& =&3\cdot2\left(\cos\left(\dfrac{\left. \dfrac{\pi}{6}t+\dfrac{\pi}{6}+\dfrac{\pi}{6}t+\dfrac{\pi}{4}\right)}{2}\right)\cos\left(\dfrac{\dfrac{\pi}{6}t+\dfrac{\pi}{6}-\dfrac{\pi}{6}t-\dfrac{\pi}{4}}{2}\right)\right)& =&6\cos\left(\dfrac{\pi}{6}t+\dfrac{5\pi}{24}\right)\cdot\cos\left(-\dfrac{\pi}{24}\right)=6\cos\left(\dfrac{\pi}{24}\right)\cdot\cos\left(\dfrac{\pi}{6}t+\dfrac{5\pi}{24}\right)$.
 \end{eqnarray*}
 Vậy pha ban đầu của dao động tổng hợp là $\dfrac{5 \pi}{24}(\mathrm{rad})=37,5^\circ$.
}
\end{ex}

% ===== Câu 9 =====
% ID: T11C1B3-04-TN
% Mức: VDC
\begin{ex}
% ID: T11C1B3-04-TN
% Mức: VDC
Nếu $\tan \dfrac{x}{2}=\dfrac{a}{b}$ thì biểu thức $a\sin x+b\cos x$ bằng
\choice
{$a$}
{\True $b$}
{$\dfrac{a+b}{a}$}
{$\dfrac{a+b}{b}$}
\loigiai{
Đặt $t=\tan \dfrac{x}{2}=\dfrac{a}{b}$ nên 
 \begin{eqnarray*}
 & $\sin$ x= $\dfrac{2t}{1+t^2}=\dfrac{2\dfrac{a}{b}}{1+\dfrac{a^2}{b^2}}=\dfrac{2ab}{a^2+b^2}$ & $\cos$ x= $\dfrac{1-t^2}{1+t^2}=\dfrac{1-\dfrac{a^2}{b^2}}{1+\dfrac{a^2}{b^2}}=\dfrac{b^2-a^2}{a^2+b^2}$.
 \end{eqnarray*}
 Vậy $a\sin x+b\cos x=\dfrac{2a^2b}{a^2+b^2}+\dfrac{b^3-a^2b}{a^2+b^2}=b$.
}
\end{ex}

% ===== Câu 10 =====
% ID: T11C1B3-05-DS
% Mức: VD
\dangbt{Tìm giá trị lớn nhất của biểu thức lượng giác}
\begin{ex}
% ID: T11C1B3-05-DS
% Mức: VD
Biến đổi được các biểu thức sau về dạng tích số. Khi đó
\choiceTF

\loigiai{
A. ch Sai. $\cos 3 x+\cos x=2 \cos 2 x \cdot \cos x$.
B. ch Sai. $\sin 3 x+\sin 2 x=2 \sin \dfrac{5 x}{2} \cos \dfrac{x}{2}$.
C. ch Đúng. $\cos 4 x-\cos x=-2 \sin \dfrac{5 x}{2} \sin \dfrac{3 x}{2}$.
D. ch Đúng. $\sin 5 x-\sin x=2 \cos 3 x \sin 2 x$.
}
\end{ex}

% ===== Câu 11 =====
% ID: T11C1B3-05-TN
% Mức: VD
\dangbt{Tính toán giá trị của tang tổng hai góc}
\begin{ex}
% ID: T11C1B3-05-TN
% Mức: VD
Nếu $\tan x=0,5$; $\sin y=\dfrac{3}{5}(0\lt y<90^\circ)$ thì $\tan (x+y)$ bằng
\choice
{\True $2$}
{$3$}
{$4$}
{$5$}
\loigiai{
$\tan x=0,5=\dfrac{1}{2}$, $\sin y=\dfrac{3}{5}(0\lt y<90^\circ)\Rightarrow \cos y=\dfrac{4}{5}\Rightarrow \tan y=\dfrac{3}{4}$.
$\tan (x+y)=\dfrac{\tan x+\tan y}{1-\tan x\cdot \tan y}=\dfrac{\dfrac{1}{2}+\dfrac{3}{4}}{1-\dfrac{1}{2}\cdot \dfrac{3}{4}}=2$.
}
\end{ex}

% ===== Câu 12 =====
% ID: T11C1B3-06-TLN
% Mức: VDC
\dangbt{Sử dụng công thức lượng giác trong tam giác}
\begin{ex}
% ID: T11C1B3-06-TLN
% Mức: VDC
Cho các góc $\alpha$, $\beta$ thỏa mãn $\dfrac{\pi}{2}\lt \alpha$, $\beta<\pi$, $\sin \alpha=\dfrac{1}{3}$, $\cos \beta=-\dfrac{2}{3}$.
Giá trị $\sin (\alpha +\beta)=-\dfrac{a+b\sqrt{c}}{d}$. Khi đó $bc-ad$ bằng bao nhiêu?
\shortans{$2$}
\loigiai{
Do $\dfrac{\pi}{2}\lt \alpha$, $\beta<\pi \Rightarrow \left\{ \begin{aligned} \cos \alpha<0 \\ \sin \beta>0. \end{aligned} \right.$
Ta có $\cos \alpha=-\sqrt{1-\sin ^2 \alpha}=-\sqrt{1-\dfrac{1}{9}}=-\dfrac{2 \sqrt{2}}{3};\ \sin \beta=\sqrt{1-\cos ^2 \beta}=\sqrt{1-\dfrac{4}{9}}=\dfrac{\sqrt{5}}{3}$.
Suy ra $\sin (\alpha+\beta)=\sin \alpha \cdot \cos \beta+\cos \alpha \cdot \sin \beta=\dfrac{1}{3} \cdot\left(-\dfrac{2}{3}\right)+\left(-\dfrac{2 \sqrt{2}}{3}\right) \cdot \dfrac{\sqrt{5}}{3}=-\dfrac{2+2 \sqrt{10}}{9}$.
Do đó $\left\{ \begin{aligned} a &=2 \\ b &=2 \\ c &=10 \\ d &=9 \end{aligned} \right.\Rightarrow bc-ad=2$.
}
\end{ex}

% ===== Câu 13 =====
% ID: T11C1B3-07-TLN
% Mức: VDC
\begin{ex}
% ID: T11C1B3-07-TLN
% Mức: VDC
Cho tam giác $ABC$ có $AB=c$, $AC=b$, $BC=a$, thỏa mãn $\dfrac{-1+\cos (\pi+B)}{\sin (\pi-A-C)}=-\dfrac{2 a+c}{\sqrt{4 a^2-c^2}}.$
 Giá trị của biểu thức $(a-b)^2c^3$ bằng
\shortans{$0$}
\loigiai{
Ta có
 \begin{eqnarray*}
 && $\dfrac{-1+\cos (\pi +B)}{\sin (\pi -A-C)}=-\dfrac{2a+c}{\sqrt{4a^2-c^2}}\Leftrightarrow\dfrac{-1-\cos B}{\sin (\pi -(A+C))}=-\dfrac{2a+c}{\sqrt{4a^2-c^2}}$ &\Leftrightarrow& $\dfrac{1+\cos B}{\sin B}=\dfrac{2a+c}{\sqrt{4a^2-c^2}}\Leftrightarrow\dfrac{1+\cos B}{\sin B}=\dfrac{2a+c}{\sqrt{4a^2-c^2}}\Leftrightarrow{\left(\dfrac{1+\cos B}{\sin B}\right)^2}=\left(\dfrac{2a+c}{\sqrt{4a^2-c^2}}\right)^2$ &\Leftrightarrow& $\dfrac{(1+\cos B)^2}{\sin^2B}=\dfrac{(2a+c)^2}{4a^2-c^2}\Leftrightarrow\dfrac{(1+\cos B)^2}{1-\cos^2B}=\dfrac{(2a+c)^2}{4a^2-c^2}\Leftrightarrow\dfrac{1+\cos B}{1-\cos B}=\dfrac{2a+c}{2a-c}$ &\Leftrightarrow& $\dfrac{1+\cos B}{1-\cos B}-1=\dfrac{2a+c}{2a-c}-1\Leftrightarrow$ 2ac $\cdot\cos$ B=c^2\Leftrightarrow a^2+c^2-b^2=c^2\Leftrightarrow a^2=b^2\Leftrightarrow a=b.
 \end{eqnarray*}
 Vậy $(a-b)^2c^3=0$.
}
\end{ex}

% ===== Câu 14 =====
% ID: T11C1B3-08-TLN
% Mức: VD
\dangbt{Tìm giá trị lớn nhất của biểu thức lượng giác}
\begin{ex}
% ID: T11C1B3-08-TLN
% Mức: VD
Cho $\alpha-\beta=\dfrac{\pi}{3}$. Giá trị của biểu thức $B=(\cos \alpha+\sin \beta)^2+(\cos \beta-\sin \alpha)^2$ bằng $a-\sqrt{b}$. Khi đó $a+b$ bằng
\shortans{$5$}
\loigiai{
Ta có
 \begin{eqnarray*}
 $B& =&\cos^2\alpha+\sin^2\beta+2\cos\alpha\sin\beta+\cos^2\beta+\sin^2\alpha-2\sin\alpha\cos\beta& =&2+2\sin\beta\cos\alpha-2\cos\beta\sin\alpha& =&2+2\sin(\beta-\alpha)=2+2\sin\left(-\dfrac{\pi}{3}\right)=2-\sqrt{3}$.
 \end{eqnarray*}
}
\end{ex}

% ===== Câu 15 =====
% ID: T11C1B3-09-TLN
% Mức: VDC
\begin{ex}
% ID: T11C1B3-09-TLN
% Mức: VDC
Biết $\dfrac{3-4\cos 2a+\cos 4a}{3+4\cos 2a+\cos 4a}=k\tan^4a$, tìm $k$?
\shortans{$1$}
\loigiai{
Ta có 
 \begin{eqnarray*}
 $\dfrac{3-4 \cos 2 a+\cos 4 a}{3+4 \cos 2 a+\cos 4 a}&=&\dfrac{2 \cos ^2 2 a-4 \cos 2 a+2}{2 \cos ^2 2 a+4 \cos 2 a+2}=\dfrac{2(\cos ^2 2 a-2 \cos 2 a+1)}{2(\cos ^2 2 a+2 \cos 2 a+1)}&=&\dfrac{(1-\cos 2 a)^2}{(1+\cos 2 a)^2}=\dfrac{\left[1-(1-2 \sin ^2 a)\right]^2}{\left[1+(2 \cos ^2 a-1)\right]^2}=\dfrac{4 \sin ^4 a}{4 \cos ^4 a}=\tan^4\,\text{a}$.
 \end{eqnarray*}
}
\end{ex}

% ===== Câu 16 =====
% ID: T11C1B3-10-TLN
% Mức: VD
\dangbt{Tính toán giá trị của tang tổng hai góc}
\begin{ex}
% ID: T11C1B3-10-TLN
% Mức: VD
Cho $\cot \alpha=3$. Tính $A=\dfrac{3 \sin \alpha-2 \cos \alpha}{12 \sin ^3 \alpha+4 \cos ^3 \alpha}$ (kết quả làm tròn đến hàng phần mười).
\shortans{$-0{,}3$}
\loigiai{
$A=\dfrac{3 \sin \alpha-2 \cos \alpha}{12 \sin ^3 \alpha+4 \cos ^3 \alpha}=\dfrac{\dfrac{1}{\sin ^2 \alpha}(3-2 \cot \alpha)}{12+4 \cot ^3 \alpha}=(1+\cot ^2 \alpha) \dfrac{3-2 \cot \alpha}{12+4 \cot ^3 \alpha}=-\dfrac{1}{4}\approx -0,3$.
}
\end{ex}

% ===== Câu 17 =====
% ID: T11C1B3-11-TLN
% Mức: VD
\begin{ex}
% ID: T11C1B3-11-TLN
% Mức: VD
Cho hai góc nhọn $a$ và $b$ với $\tan a=\dfrac{1}{7}$ và $\tan b=\dfrac{3}{4}$. Tính $a+b$.
\shortans{$45$}
\loigiai{
Ta có $\tan (a+b)=\dfrac{\tan a+\tan b}{1-\tan a \tan b}=\dfrac{\dfrac{1}{7}+\dfrac{3}{4}}{1-\dfrac{1}{7} \cdot \dfrac{3}{4}}=1$, suy ra $a+b=45^\circ$.
}
\end{ex}
